\section{Catálogo de requisitos candidatos}
\label{sec:catalogo-requisitos}

El caso de estudio corresponde a los procesos de matrícula y registro de calificaciones de una unidad educativa. Se contemplan estudiantes, representantes, docentes, personal de secretaría, coordinación académica y autoridades como partes interesadas en los procesos escolares. Los requisitos se especifican como comportamientos y condiciones observables mediante el verbo modal \emph{deberá}, de conformidad con el enfoque de especificación individual \cite{iso29148,wiegers2013}.

Se analizan doce requisitos funcionales (RF) y ocho requisitos de calidad (RNF). Las cifras utilizadas en los RNF son \emph{umbrales propuestos}, no mediciones de una unidad educativa específica. El alcance está limitado a las funciones académicas de matrícula y calificaciones.

\begingroup
\footnotesize
\renewcommand{\arraystretch}{1.12}
\begin{longtable}{@{}p{0.12\linewidth}p{0.83\linewidth}@{}}
\caption{Catálogo de veinte requisitos candidatos para la unidad educativa}\label{tab:catalogo-veinte}\\
\toprule
\textbf{ID} & \textbf{Requisito propuesto} \\
\midrule
\endfirsthead
\multicolumn{2}{@{}l}{\textit{Continuación del catálogo de requisitos}}\\
\toprule
\textbf{ID} & \textbf{Requisito propuesto} \\
\midrule
\endhead
\bottomrule
\endlastfoot
\textbf{RF-01} & El sistema deberá registrar un estudiante con identificador institucional único cuando la secretaría confirme los datos de matrícula y la identificación no exista, con el fin de crear su expediente académico. \\
\textbf{RF-02} & El sistema deberá autenticar a un usuario mediante credenciales institucionales cuando solicite ingresar a la plataforma, con el fin de habilitar las funciones correspondientes a su perfil. \\
\textbf{RF-03} & El sistema deberá mostrar los paralelos ofertados en el período lectivo seleccionado cuando un usuario autorizado consulte la oferta académica, con el fin de identificar opciones de matrícula. \\
\textbf{RF-04} & El sistema deberá registrar la matrícula de un estudiante en el paralelo seleccionado cuando exista un cupo disponible y el estudiante cumpla las condiciones académicas definidas, con el fin de formalizar su inscripción. \\
\textbf{RF-05} & El sistema deberá rechazar una solicitud de matrícula en un paralelo cuando el número de matriculados alcance su capacidad definida, con el fin de evitar sobreasignación de cupos. \\
\textbf{RF-06} & El sistema deberá generar un comprobante con código, estudiante, paralelo y período lectivo cuando la matrícula quede registrada, con el fin de acreditar la inscripción. \\
\textbf{RF-07} & El sistema deberá asociar un docente autorizado a una asignatura de un paralelo cuando la coordinación confirme su distribución académica, con el fin de habilitar el registro de notas. \\
\textbf{RF-08} & El sistema deberá registrar la calificación numérica de una evaluación para un estudiante matriculado cuando el docente asignado confirme el valor dentro de la escala institucional configurada, con el fin de mantener el registro académico. \\
\textbf{RF-09} & El sistema deberá registrar una nueva versión de una calificación con la justificación de modificación cuando el responsable autorizado apruebe su corrección, con el fin de conservar el historial de cambios. \\
\textbf{RF-10} & El sistema deberá mostrar al estudiante autenticado sus calificaciones publicadas del período seleccionado cuando consulte su historial, con el fin de conocer su desempeño académico. \\
\textbf{RF-11} & El sistema deberá generar el boletín de un estudiante con asignaturas y calificaciones publicadas cuando secretaría solicite el documento para un período lectivo, con el fin de respaldar su trayectoria académica. \\
\textbf{RF-12} & El sistema deberá generar un reporte de matriculados por paralelo y período cuando la dirección académica seleccione un período lectivo, con el fin de supervisar la ocupación de cupos. \\
\textbf{RNF-01} & El sistema deberá completar al menos el 95 \% de los registros de matrícula en un máximo de 3 segundos cuando existan 100 solicitudes simultáneas sobre una base de 3000 estudiantes, con el fin de mantener tiempos controlados en el inicio de período. \\
\textbf{RNF-02} & El sistema deberá mantener una disponibilidad mensual mínima del 99,5 \% del portal cuando se evalúe un mes calendario de operación, incluyendo mantenimientos, con el fin de asegurar acceso regular a servicios académicos. \\
\textbf{RNF-03} & El sistema deberá impedir el acceso a calificaciones individuales de un estudiante desde cuentas no autorizadas en el 100 \% de los casos cuando se ejecute una batería de 50 solicitudes negativas con perfiles de estudiante, docente y secretaría, con el fin de proteger información académica personal. \\
\textbf{RNF-04} & El sistema deberá permitir que al menos el 90 \% de una muestra de 20 usuarios adultos complete una matrícula simulada en no más de 5 minutos cuando utilicen por primera vez la interfaz sin ayuda, con el fin de favorecer la autogestión. \\
\textbf{RNF-05} & El sistema deberá mostrar sin desplazamiento horizontal las pantallas de consulta de matrícula y notas en al menos cinco resoluciones de pantalla entre 360 y 768 píxeles cuando se pruebe en navegadores modernos, con el fin de facilitar su uso en teléfonos. \\
\textbf{RNF-06} & El sistema deberá generar al menos una copia íntegra de las bases de matrícula y calificaciones cada 24 horas cuando se encuentre en operación, con el fin de limitar la pérdida de datos ante fallos. \\
\textbf{RNF-07} & El sistema deberá restaurar el servicio de consulta de calificaciones en un máximo de 60 minutos cuando se simule una caída completa del servidor en un ambiente de recuperación, con el fin de reducir la interrupción académica. \\
\textbf{RNF-08} & El sistema deberá registrar el 100 \% de las modificaciones de calificaciones con identificador de actor, momento, valor anterior y valor nuevo cuando se ejecuten 50 cambios autorizados de prueba, con el fin de permitir auditorías internas. \\
\end{longtable}
\endgroup

\subsection{Criterios de revisión de los requisitos}

La revisión de cada requisito considera el comportamiento solicitado, su condición de activación, el propósito, la autorización del actor y la posibilidad de comprobarlo mediante pruebas. En los requisitos de calidad se debe verificar explícitamente la métrica, el umbral, la unidad y el escenario de medición. El archivo \texttt{datos/requisitos.csv} contiene el código, tipo, nombre, redacción y forma propuesta de verificación, lo cual permite cotejar este catálogo con los datos empleados en las técnicas de priorización.

Las políticas de matrícula, promoción, escalas de calificaciones, cupos y roles de acceso deberán contrastarse con la normativa y la organización de la unidad educativa donde se aplique el ejercicio. Esta necesidad de validación no implica que se hayan efectuado consultas institucionales.
